\section{Candidate Physical System and Experimental Test Strategy}
\label{sec:candidate_system}

The general framework becomes substantially more concrete when applied to known nuclear isomers. Among presently known systems, the low-lying isomer of $^{235}$U, denoted $^{235m}$U, provides a particularly useful test case because both ingredients required by the proposed mechanism are already partially established: the isomer possesses a fission response different from that of the nuclear ground state, and mechanisms for externally populating the isomer have been experimentally investigated or theoretically proposed.

Importantly, the use of $^{235m}$U in connection with fission control is not itself a new suggestion. Previous work has explicitly noted that the isomer may be exploited to control nuclear fission because its thermal-neutron fission cross section differs substantially from that of ground-state $^{235}$U \cite{DEer1988,LiuWang2022}. The purpose of the present framework is therefore more specific: to treat the externally controllable nuclear-state population as a dynamical control variable, connect it explicitly to an effective fission cross section, and propagate that microscopic modification into macroscopic quantities such as fission rate, neutron multiplication, and reactor reactivity.

\subsection{\texorpdfstring{$^{235m}$U}{235mU} as a candidate system}

The first excited nuclear state of $^{235}$U lies only
\begin{equation}
E_m
=
76.737\pm0.018\ \mathrm{eV}
\end{equation}
above the ground state \cite{Ponce2018}. This exceptionally low excitation energy is unusual for a nuclear transition. The state has a half-life of approximately 26 min and decays predominantly through internal conversion.

For the present proposal, however, the most important property is its different neutron-induced fission response.

D'Eer \textit{et al.} measured the fission cross section of $^{235m}$U relative to that of ground-state $^{235}$U and obtained
\begin{equation}
\frac{\sigma_f^m}{\sigma_f^g}
=
1.61\pm0.44
\end{equation}
for cold neutrons and
\begin{equation}
\frac{\sigma_f^m}{\sigma_f^g}
=
2.47\pm0.45
\end{equation}
for thermal neutrons \cite{DEer1988}.

Thus,
\begin{equation}
\boxed{
\sigma_f^m\neq\sigma_f^g
}
\end{equation}
is experimentally established.

This result demonstrates that the neutron-induced fission response of a fissile isotope is not determined by isotopic identity alone; it can depend strongly on the internal nuclear state.

\subsection{Minimal two-population description}

Let $p_m(t)$ be the fraction of $^{235}$U nuclei occupying the isomeric state and
\begin{equation}
p_g(t)=1-p_m(t)
\end{equation}
the fraction in the ground state.

The effective microscopic fission cross section is
\begin{equation}
\sigma_f^{\mathrm{eff}}(E,t)
=
[1-p_m(t)]\,\sigma_f^g(E)
+
p_m(t)\,\sigma_f^m(E).
\end{equation}

Equivalently,
\begin{equation}
\sigma_f^{\mathrm{eff}}
=
\sigma_f^g
+
p_m
\left(
\sigma_f^m-\sigma_f^g
\right).
\end{equation}

Defining
\begin{equation}
r(E)
=
\frac{\sigma_f^m(E)}{\sigma_f^g(E)},
\end{equation}
one obtains
\begin{equation}
\frac{\Delta\sigma_f}{\sigma_f^g}
=
p_m[r(E)-1].
\end{equation}

For thermal neutrons, using the measured value $r\simeq2.47$,
\begin{equation}
\boxed{
\frac{\Delta\sigma_f}{\sigma_f^g}
\simeq
1.47\,p_m
}.
\label{eq:u235_thermal_effect}
\end{equation}

Equation~\eqref{eq:u235_thermal_effect} exposes the control problem directly. Once the state-dependent fission cross sections are known, the new dynamical quantity is the controllable isomeric fraction $p_m(t)$:
\begin{equation}
u_Q(t)
\longrightarrow
p_m(t)
\longrightarrow
\sigma_f^{\mathrm{eff}}(t).
\end{equation}

The key microscopic feasibility question becomes whether $p_m(t)$ can be externally and reproducibly controlled.

\subsection{External population of the isomer}

Producing a useful controlled population of $^{235m}$U remains the principal unresolved step.

Nuclear excitation by electronic transition (NEET) has been investigated as one possible mechanism. A dedicated search using a laser-generated uranium plasma observed no statistically significant $^{235m}$U production attributable to NEET and established an upper bound under the experimental conditions \cite{Chodash2016}.

Liu and Wang later investigated direct inelastic scattering of low-energy electrons using a distorted-wave treatment \cite{LiuWang2022}. Their calculations predict excitation cross sections larger than earlier plane-wave estimates and also discuss indirect population through higher nuclear levels as a possible route.

These studies do not establish a practical actuator. They instead define the open problem: whether an externally adjustable interaction can generate a sufficiently large, reproducible, and sufficiently rapid modulation of the isomeric fraction.

The present framework does not require a particular excitation technology. What matters is a reproducible mapping
\begin{equation}
u_Q(t)
\rightarrow
p_m(t).
\end{equation}

\subsection{Distinguishing excitation from control}

Producing an isomer is not equivalent to controlling fission through the isomeric population.

The existence of
\begin{equation}
^{235}\mathrm{U}_g
\rightarrow
^{235m}\mathrm{U}
\end{equation}
is not by itself sufficient.

A control mechanism requires
\begin{equation}
p_m=p_m(u_Q)
\end{equation}
and a corresponding modification of neutron-induced fission under an otherwise fixed neutron field:
\begin{equation}
\boxed{
u_Q
\rightarrow
p_m
\rightarrow
\sigma_f^{\mathrm{eff}}
\rightarrow
R_f
}.
\end{equation}

This distinction is what connects nuclear excitation physics to control theory.

\subsection{Minimal falsification experiment}

A decisive test does not require a multiplying reactor.

Consider two otherwise equivalent fissile ensembles exposed to the same neutron spectrum,
\begin{equation}
\Phi_1(E)=\Phi_2(E),
\end{equation}
but with different isomeric fractions,
\begin{equation}
p_m^{(1)}
\neq
p_m^{(2)}.
\end{equation}

The predicted fission-rate difference is
\begin{equation}
\Delta R_f
=
N_f
\int
\Phi(E)
\left[
p_m^{(2)}-p_m^{(1)}
\right]
\left[
\sigma_f^m(E)-\sigma_f^g(E)
\right]
\,dE.
\label{eq:u235_experiment}
\end{equation}

Since the cross-section difference is already known to be nonzero for cold and thermal neutrons, the decisive experimental step is to create
\begin{equation}
\Delta p_m
=
p_m^{(2)}-p_m^{(1)}
\end{equation}
through an externally adjustable actuator.

The hypothesis predicts
\begin{equation}
\boxed{
\left.
\frac{\partial R_f}
{\partial u_Q}
\right|_{\Phi}
\neq0
}
\end{equation}
provided that
\begin{equation}
\frac{\partial p_m}{\partial u_Q}\neq0
\end{equation}
and
\begin{equation}
\sigma_f^m-\sigma_f^g\neq0.
\end{equation}

Observation of a reproducible chain
\begin{equation}
u_Q
\rightarrow
p_m
\rightarrow
R_f
\end{equation}
with the neutron spectrum held fixed would constitute direct evidence of externally actuated nuclear-state control of neutron-induced fission.

\subsection{From microscopic control to macroscopic relevance}

A positive microscopic result would not imply practical reactor control.

Let $f_Q$ denote the fraction of fissile nuclei whose state population can effectively be modified. To first order,
\begin{equation}
\frac{\delta\Sigma_f}{\Sigma_f}
\sim
f_Q
\frac{\delta\sigma_f}{\sigma_f}.
\end{equation}

Macroscopic relevance therefore depends not only on
\begin{equation}
\Delta\sigma_f
=
\sigma_f^m-\sigma_f^g,
\end{equation}
but also on the achievable state population, actuation and relaxation timescales, spatial penetration, and energetic cost.

A minimal feasibility parameter set is
\begin{equation}
\boxed{
\left\{
p_m^{\max},
\frac{dp_m}{dt},
\tau_m,
\Delta\sigma_f(E),
E_Q,
f_Q
\right\}
}.
\end{equation}

Only once these quantities are constrained does it become meaningful to propagate the microscopic effect into $\delta k_{\mathrm{eff}}$ and $\rho_Q$.

\subsection{Scope of the present contribution}

The $^{235m}$U literature requires a careful novelty statement.

Previous studies established that the isomer has a different neutron-induced fission cross section, and the possibility of exploiting this fact for fission control has been explicitly mentioned \cite{DEer1988,LiuWang2022}. The present work therefore does \emph{not} claim the first proposal to use $^{235m}$U to influence fission.

The distinct objective is to formulate nuclear-state population as an explicit dynamical control variable and embed it in the multiscale chain
\begin{equation}
u_Q(t)
\rightarrow
\rho_N(t)
\rightarrow
\sigma_f^{\mathrm{eff}}(E,t)
\rightarrow
R_f(t)
\rightarrow
k_{\mathrm{eff}}(t)
\rightarrow
\rho_Q(t).
\label{eq:multiscale_chain}
\end{equation}

Within this framework, an isomer is not merely a state with unusual reaction properties. Its population becomes a candidate \emph{state variable of the controlled fission system}.

This formulation separates three questions:
\begin{enumerate}
\item \textbf{State dependence:} do different nuclear states have different fission responses?
\item \textbf{Actuation:} can their populations be externally and dynamically controlled?
\item \textbf{System coupling:} can the resulting microscopic modification influence macroscopic neutron multiplication?
\end{enumerate}

For $^{235m}$U, the first question already has a positive experimental answer. The second remains an active nuclear-physics problem. The third is the reactor-physics and control-theory problem formulated here.

\subsection{Beyond \texorpdfstring{$^{235m}$U}{235mU}}

A suitable candidate system should combine two properties:
\begin{equation}
\boxed{
\text{nuclear-state controllability}
\times
\text{fission sensitivity to that state}
}.
\end{equation}

Possible candidate classes include long-lived spectroscopic isomers, shape isomers, and selected excited configurations of actinides. A state with a large fission effect but negligible controllability is unsuitable; a readily controllable state is equally unhelpful if its fission response is unchanged.

\subsection{Experimental decision tree}

The research program can proceed through increasingly stringent tests:
\begin{enumerate}
\item demonstrate reproducible external control of a fission-distinct nuclear-state population,
\begin{equation}
u_Q\rightarrow\rho_N;
\end{equation}
\item under a fixed neutron spectrum, verify that the controlled state changes neutron-induced fission,
\begin{equation}
\rho_N
\rightarrow
\sigma_f^{\mathrm{eff}}
\rightarrow
R_f;
\end{equation}
\item determine whether the achievable controlled fraction is sufficient to perturb neutron multiplication,
\begin{equation}
R_f
\rightarrow
\delta k_{\mathrm{eff}}.
\end{equation}
\end{enumerate}

This hierarchy separates the fundamental microscopic question from the engineering one. A negative result at the first or second stage would constrain or falsify the proposed mechanism without reactor-scale experiments. A positive result would provide the parameters needed to assess macroscopic relevance.
